% Vista científica exacta materializada desde una rama condicional.
% Fuente: source/demostraciones_primarias/31_06c_sectores_topologicos_no_omision.tex; rama: HMTIncludeTopologicalBlock.
% No modifica el contenido matemático de la rama de origen.
\chapter[Incidence and Topological Sectors]{Incidence, Fibonacci and topological sectors}
\label{chap:incidencia-fibonacci-topologia}
\index{topological sector}\index{Fibonacci, ring based}
\index{braid group}

The residue--support incidence prolongs the atlas toward fusion and braid
structures. The order is essential: first the based Fibonacci ring
and its dimensions are constructed; the algebraic representations
of braids and the pointed center are then studied. Each object preserves its own
domain and is not identified with a physical particle merely because it shares a name or
cardinal.

\section{Residue-incident support}
\label{sec:incidencia-residuo-soporte}

We fix the residue classes and the supports in the order
\[
 \begin{aligned}
 C_1&=\{1,4,7\},&
 C_2&=\{2,5,8\},&
 C_0&=\{3,6,9\},\\
 \mathcal D_{\rm act}&=\{1,3,5,7\},&
 \mathcal D_{\rm evt}&=\{2,4,6,8\},&
 \mathcal D_{\bar0}^{(9)}&=\{9\}.
 \end{aligned}
\]
The final symbol means ‘support of the null nonadic residue’. It does not designate
the complete domain \(D_9=\{1,\ldots,9\}\) or a physical zero-mass sector.

Let \(E_{\mathcal D}\cong\mathbb Z^3\) and
\(E_C\cong\mathbb Z^3\) be the free modules with ordered bases
\[
 (\mathcal D_{\rm act},\mathcal D_{\rm evt},
   \mathcal D_{\bar0}^{(9)})
 \quad\text{and}\quad
 (C_1,C_2,C_0),
\]
respectively. The incidence map
\(B_{\rm res,sup}:E_{\mathcal D}\to E_C\) counts intersections.

\begin{theorem}[Smith form and image of incidence]
\label{thm:incidencia-residuo-soporte}
In the preceding bases,
\[
 B_{\rm res,sup}
 =\bigl(|C_i\cap\mathcal D_j|\bigr)_{i,j}
 =
 \begin{pmatrix}
 2&1&0\\
 1&2&0\\
 1&1&1
 \end{pmatrix}.
\]
Moreover,
\[
 \det B_{\rm res,sup}=3,\qquad
 \operatorname{SNF}(B_{\rm res,sup})
 =\operatorname{diag}(1,1,3),
\]
\[
 \operatorname{im}B_{\rm res,sup}
 =\{(u,v,w)\in\mathbb Z^3:u+v\equiv0\pmod3\}.
\]
If
\[
 P=\begin{pmatrix}0&1&0\\1&0&0\\0&0&1\end{pmatrix},
\]
the simultaneous exchange
\(C_1\leftrightarrow C_2\) and
\(\mathcal D_{\rm act}\leftrightarrow\mathcal D_{\rm evt}\) is equivariant:
\[
 P B_{\rm res,sup}P=B_{\rm res,sup}
 \qquad\Longleftrightarrow\qquad
 PB_{\rm res,sup}=B_{\rm res,sup}P.
\]
\end{theorem}

\begin{proof}
The nine entries follow from the intersections.
\[
\begin{array}{cccc}
 &\mathcal D_{\rm act}&\mathcal D_{\rm evt}
 &\mathcal D_{\bar0}^{(9)}\\
 C_1&\{1,7\}&\{4\}&\varnothing\\
 C_2&\{5\}&\{2,8\}&\varnothing\\
 C_0&\{3\}&\{6\}&\{9\}.
\end{array}
\]
The expansion by the third column gives
\[
 \det B_{\rm res,sup}
 =\det\begin{pmatrix}2&1\\1&2\end{pmatrix}=3.
\]
The greatest common divisor of the entries is one. The smallest second-order
minor formed by the first and third rows and the second and third columns
has determinant one:
\[
 \det\begin{pmatrix}1&0\\1&1\end{pmatrix}=1.
\]
By the Smith divisors, \(d_1=d_2=1\). The product
\(d_1d_2d_3\) agrees with the absolute value of the determinant, which the
preceding calculation fixes at three. Consequently,

\[
 d_1d_2d_3=3.
\]
The normal form is therefore \(\operatorname{diag}(1,1,3)\).

Each column satisfies \(u+v\equiv0\pmod3\), so the image is contained
in the indicated submodule. Conversely, if \(u+v\equiv0\pmod3\), then
\[
 x=\frac{2u-v}{3},\qquad
 y=\frac{2v-u}{3},\qquad
 z=w-x-y
\]
are integers and
\[
 B_{\rm res,sup}(x,y,z)^{\mathsf T}=(u,v,w)^{\mathsf T}.
\]
This proves equality of images. Finally, explicit multiplication
by \(P\) exchanges the first two rows and the first two columns and
leaves the matrix invariant.
\end{proof}

\begin{criterion}[Incidence control]
\label{crit:incidencia-residuo-soporte}
The theorem would be falsified by a single miscounted intersection, a minor that altered
the Smith divisors, a vector of the image with \(u+v\not\equiv0\pmod3\),
a vector satisfying the congruence but admitting no exhibited preimage,
or the loss of equivariance \(PBP=B\).
\end{criterion}

\section[Fibonacci-based Ring]{Based ring associated with the incidence
  \texorpdfstring{\(F_{\rm av}\)}{F\_av}}
\label{sec:anillo-basado-fibonacci}

The \cref{thm:descenso-necesario-fav} proves, from the orbit of
modes \((\Sigma,\Pi,\Pi)\), the matrix
\[
 F_{\rm av}=
 \begin{pmatrix}0&1\\1&1\end{pmatrix}.
\]
From this matrix the following based ring is obtained.

\begin{definition}[Fibonacci-based Ring]
\label{def:anillo-basado-fibonacci}
Let
\[
 K_{\rm Fib}:=\mathbb Z\mathbf 1\oplus\mathbb Z\tau
\]
the free abelian group with unit \(\mathbf 1\) and bilinear product fixed by
\[
 \tau\cdot\mathbf 1=\tau,\qquad
 \tau\cdot\tau=\mathbf 1+\tau.
\]
In this stratum, \(+\) and \(\cdot\) are the addition and multiplication of the
ring. The categorical symbols \(\oplus\) and \(\otimes\) are reserved for
a categorification chosen subsequently.
\end{definition}

\begin{theorem}[Realization of \(F_{\rm av}\) and Perron dimension]
\label{thm:anillo-fibonacci-fav}
In the ordered basis \((\mathbf 1,\tau)\), multiplication by \(\tau\) in
\(K_{\rm Fib}\) has matrix \(F_{\rm av}\). If
\(F_0=0\), \(F_1=1\), and
\(F_{n+1}=F_n+F_{n-1}\) for \(n\geq1\), then
\[
 \tau^n=F_{n-1}\mathbf 1+F_n\tau
 \qquad(n\geq1).
\]
The unique homomorfismo positivo of dimension that envia
\(\mathbf 1\) to one satisfies
\[
 \operatorname{FPdim}(\tau)=\varphi.
\]
Identification of the two typed sheets with the basis
\((\mathbf 1,\tau)\) completely determines this ring and its positive
dimension.
\end{theorem}

\begin{proof}
The columns of the matrix of multiplication by \(\tau\) are the coordinates
of
\[
 \tau\cdot\mathbf 1=\tau
 \quad\text{and}\quad
 \tau\cdot\tau=\mathbf 1+\tau;
\]
they are therefore \((0,1)^{\mathsf T}\) and \((1,1)^{\mathsf T}\), respectively.
Thus one obtains \(F_{\rm av}\).

The power formula holds for \(n=1\). If it holds for \(n\), then
\[
 \tau^{n+1}
 =F_{n-1}\tau+F_n(\mathbf 1+\tau)
 =F_n\mathbf 1+F_{n+1}\tau,
\]
which closes the induction. Finally, a positive multiplicative dimension
\(d=\operatorname{FPdim}(\tau)\) obeys
\[
 d^2=1+d.
\]
Its roots are \(\varphi\) and \(-\varphi^{-1}\); positivity selects
\(d=\varphi\).
\end{proof}

The designation \emph{Fibonacci} is thus explained by the fusion
rule itself: upon iterating \(\tau\), the two possible channels appear with
consecutive multiplicities \(F_{n-1}\) and \(F_n\). The space of fusion states
therefore grows with asymptotic ratio \(\varphi\). This is the
combinatorial basis of its computational use: information is encoded in global
fusion channels and not in a mass assigned to a point.

\begin{criterion}[Conditions for a Fibonacci categorification]
\label{crit:categorificacion-fibonacci}
The annular layer would be refuted if the multiplication matrix were not
\(F_{\rm av}\), if the power recurrence failed, or if the positive
dimension were not \(\varphi\). A braided categorification whose
Grothendieck ring is that of \cref{def:anillo-basado-fibonacci} additionally requires a
monoidal functor
\[
 \mathfrak C_{\rm Fib}:\mathsf{TPK}
 \dashrightarrow\mathcal C_{\rm Fib},
\]
the verification of the pentagon and hexagon in its image and an internal rule that
selects gauge and chirality.
These conditions are not consequences of the annular identity by itself.
\end{criterion}

\section[Braid Characters]{Six characters and four braid twists}
\label{sec:rescate-trenzas-ex02}

Let \(B_n\), \(n\geq2\), be the Artin braid group with generators
\(\sigma_1,\ldots,\sigma_{n-1}\) and relations
\[
 \sigma_i\sigma_{i+1}\sigma_i
 =\sigma_{i+1}\sigma_i\sigma_{i+1},
 \qquad
 \sigma_i\sigma_j=\sigma_j\sigma_i\quad(|i-j|\geq2),
\]
and let \(w:B_n\to\mathbb Z\) be the sum of the exponents. The readers
\[
 w_2:=w\bmod2,\qquad w_3:=w\bmod3
\]
are brought together, by the Chinese remainder theorem, in
\[
 w_6:B_n\longrightarrow
 \mathbb Z/2\mathbb Z\times\mathbb Z/3\mathbb Z
 \simeq\mathbb Z/6\mathbb Z.
\]

\begin{theorem}[Abelianization, six characters and permutation torsions]
\label{thm:trenzas-caracteres-z6}
For \(n\geq2\),
\[
 B_n^{\rm ab}\simeq\mathbb Z,
\]
through \(w\). Consequently, \(w_6\) produces six unitary characters
\[
 \chi_k(b)
 :=\exp\!\left(\frac{2\pi\mathrm i k}{6}w(b)\right),
 \qquad k\in\mathbb Z/6\mathbb Z.
\]

Let \(V\ne0\) be a Hilbert space, let
\(P_i:V^{\otimes n}\to V^{\otimes n}\) be the exchange of the factors
\(i,i+1\), and let \(q\in\mu_6\). The rule
\[
 \rho_q(\sigma_i):=qP_i,
 \qquad
 \rho_q(b)=q^{w(b)}\pi_n(b),
\]
where \(\pi_n\) is the permutation representation, defines a
representation of \(B_n\). This representation descends to the symmetric group
if and only if
\[
 q^2=1.
\]
For the four values \(q\in\mu_6\setminus\{\pm1\}\),
\[
 \operatorname{Inv}(\rho_q)=0.
\]
Those four values define genuinely braided scalar torsions, but
do not by themselves constitute four Fock spaces nor four physical
anyonic sectors.
\end{theorem}

\begin{proof}
In every abelian quotient, the braid relation gives
\[
 2[\sigma_i]+[\sigma_{i+1}]
 =[\sigma_i]+2[\sigma_{i+1}],
\]
and thus all generators have the same class. No other relation remains on
that class, since \(w(\sigma_i)=1\); hence the abelianization is
\(\mathbb Z\).

The operators \(P_i\) satisfy the braid relations and
\(P_i^2=I\). Multiplying each generator by the same scalar \(q\) preserves
both relations, so \(\rho_q\) is a representation. To
factor through \(S_n\), it must satisfy the additional relation
\(\sigma_i^2=1\), and
\[
 \rho_q(\sigma_i)^2=q^2I.
\]
This proves criterion \(q^2=1\).

If \(v\) were invariant, then
\(qP_i v=v\), that is, \(P_i v=q^{-1}v\), for every \(i\). Since each
involution \(P_i\) has only eigenvalues \(\pm1\), no such nonzero vector
exists when \(q\ne\pm1\). Consequently,
\(\operatorname{Inv}(\rho_q)=0\) for the remaining four values.
\end{proof}

When \(\dim V>1\) and \(n\geq3\), the operators \(P_1\) and \(P_2\) do not
commute; hence the full representation does not become Abelian merely because
it arises from a scalar character. The character \(\chi_k\) and the representation
\(\rho_q\) are distinct objects: the former is one-dimensional; the latter
retains the permutation action.

\section[Modular center and topological realizations]{Modular center pointed and Fibonacci realization conditions Ising}
\label{sec:centro-apuntado-no-go}

Let
\[
 A=(\mathbb Z/9\mathbb Z)^2
\]
and let \(\operatorname{Vec}_A\) be the category of
\(A\)-graded vector spaces with trivial associator. Fixing a nondegenerate bicharacter
identifies \(\widehat A\) with \(A\), but does not alter the following counts.

\begin{theorem}[Pointed center and dimensional obstruction]
\label{thm:centro-apuntado-no-go}
The Drinfeld center
\[
 \mathcal Z(\operatorname{Vec}_A)
\]
is a pointed modular category whose simple objects are parameterized by
\((a,\chi)\in A\times\widehat A\). It therefore has
\[
 |A|\,|\widehat A|=81^2=6561
\]
simple, all of quantum dimension one, and total quantum dimension
\[
 \mathcal D
 =\sqrt{\sum_{x\ {\rm simple}}d_x^2}
 =\sqrt{6561}=81.
\]
In this category, neither the Fibonacci ring nor the Ising sector can be
realized while preserving simple objects and Frobenius--Perron dimensions: they require,
respectively, a simple object of dimension \(\varphi\) and one of dimension
\(\sqrt2\).
\end{theorem}

\begin{proof}
For a finite abelian group and trivial associator, every simple object of the
center is determined by a degree \(a\in A\) and a character
\(\chi\in\widehat A\) that specifies its half-braiding. Both the graded
object and the character are one-dimensional; consequently, all the
simple objects are invertible and have quantum dimension one. Since
\(|A|=|\widehat A|=9^2=81\), the number of pairs is \(81^2\), and the formula of
\(\mathcal D\) gives \(81\).

Every fusion subcategory of a pointed category is again pointed;
in particular, its simple objects have dimension one. This contradicts
\(\operatorname{FPdim}(\tau)=\varphi\) for Fibonacci and
\(d_\sigma=\sqrt2\) for Ising. The obstruction is dimensional and does not depend
on a similarity of names or cardinalities.
\end{proof}

\begin{criterion}[Conditions for a physical realization of the sectors]
\label{crit:promocion-sectores-topologicos}
The results of
\cref{thm:trenzas-caracteres-z6,thm:centro-apuntado-no-go} would be
algebraically falsified by failure of the Artin relations, the
criterion \(q^2=1\), the census \(6561\), pointedness, or
\(\mathcal D=81\). A physical realization additionally requires a
Hamiltonian or driven protocol, a spectral gap, realizable braid
operators, protection against perturbations, and a map from the
enriched multisections. None of these data follows solely from
the two preceding algebraic theorems.
\end{criterion}
